% ============================================================
\chapter{Conclusions and Future Work}
\label{chap:conclusion}
% ============================================================

\section{Summary of findings}
\label{sec:summary}

This thesis asked one question: when the correlation matrix in
hierarchical portfolio allocation is replaced by a discovered causal graph,
how much of the performance gain is attributable to the graph's skeleton
versus its edge orientations? It made the question answerable by building the
machinery it presupposes --- a reproducible rolling discovery pipeline that
delivers a byte-identical asset--asset DAG to every allocator (A1), and an
orientation-aware allocator family with provably orientation-blind skeleton
and correlation controls (A2, discharging requirement R2) --- and answered
it under multiplicity-, distribution- and stability-aware inference (A3).

The answer, in two lines: \emph{at one-year estimation windows the gain is
almost entirely the skeleton's; at two-year windows the skeleton's gain
evaporates and orientation restores it. Orientation buys robustness to the
estimation window, not level --- and no step of the decomposition is
family-wise significant against honest controls.} In detail (net annualised
Sharpe, 2007--2024, monthly rebalancing, 5\,bps costs):

\begin{enumerate}[itemsep=2pt]
  \item \textbf{At the one-year window the gain is the skeleton's.} Replacing
        the correlation matrix with the DYNOTEARS graph is worth $+0.031$ in
        total (D1 $-$ \textsc{corr}, $p=0.14$), of which $+0.022$ comes from
        the skeleton ($p=0.14$) and only $+0.009$ from orientation
        ($p=0.47$).
  \item \textbf{At the two-year window the split inverts.} The skeleton's
        contribution collapses to $+0.009$; orientation restores
        $+0.023$--$0.027$ (pairwise $p=0.05$--$0.07$, all five DYNOTEARS
        orientation contrasts positive). Orientation buys \emph{robustness to
        the estimation window}, not level: D1 ($0.411/0.395$) and D2s
        ($0.406/0.399$) are the only allocators strong at both windows, the
        increment is invariant to costs up to 20\,bps and to the graph's
        sparsity threshold, and D1 attains the highest Deflated Sharpe ratio
        of all \rsNtrials\ configurations evaluated ($\rsDoneDSR$).
  \item \textbf{Nothing survives family-wise correction against honest
        anchors.} Of the four pre-specified SPA tests, Hansen's procedure
        rejects in exactly one --- the causal family against the weak HSP
        baseline ($p=\rsSpaConsistent$); against the \textsc{corr} control
        ($p=\rsSpaSkelWtwo$/$\rsSpaSkelWfive$) and against the skeleton
        control ($p=\rsSpaDirWtwo$/$\rsSpaDirWfive$) the decomposition's
        steps are directionally consistent but statistically unconfirmed.
        The $90\%$ Model Confidence Set over \rsMcsUniverse\ one-year
        strategies eliminates only the HSP baseline.
  \item \textbf{The effect is method- and regime-specific.} Under VARLiNGAM
        graphs orientation adds nothing anywhere; under cheap ridge-Granger
        graphs the topological \emph{ordering} survives but the structural
        covariance collapses --- edge magnitudes, not just directions, are
        what expensive discovery buys. The orientation edge concentrates in
        calm (bottom-VIX-quintile) regimes; the stress hypothesis is
        rejected.
  \item \textbf{The incumbent driver route is dominated.} HSP-style driver
        selection is $K$-fragile; its monthly feedback loop is exactly inert
        across the full hyperparameter grid, and could not learn anyway
        (reward signal-to-noise ${\approx}\,\rsRewardSNR$ --- the update
        frequency \emph{is} the measurement problem); and its neural
        sensitivity step makes its Sharpe seed-dependent across
        $[\rsSeedMin,\rsSeedMax]$ over \rsSeedN\ seeds --- a range twice
        its claimed edge, with the committed value at the $\rsSeedPct$th
        percentile of its own seed distribution. The graph-based allocators
        are deterministic, and the headline three (D0, D1, D2s) exceed even
        its luckiest seed.
\end{enumerate}

\section{Contributions}
\label{sec:contributions}

\begin{itemize}[itemsep=2pt]
  \item \textbf{Methodological:} the orientation-aware allocator family ---
        structural (SEM-implied, de-standardised, residual-calibrated)
        covariance allocation; causal-ordered bisection, which replaces the
        clustering dendrogram with the DAG's topological order;
        structural-shock risk parity; co-ancestry clustering --- together
        with the fixed-graph ablation design that makes skeleton and
        orientation separately measurable, with provably orientation-blind
        controls. To my knowledge, both the family and the decomposition are
        new; the closest prior work \cite{howard2025causal} never sets
        weights from its graphs.
  \item \textbf{Empirical:} the decomposition itself (skeleton-carried at
        short windows, orientation-as-robustness at long windows, nothing
        family-wise significant against honest controls), its method- and
        regime-dependence, and the finding that textbook HRP beats the HSP
        machinery built to improve on it.
  \item \textbf{Meta-scientific:} a quantified seed audit substantiating the
        instability critique of neural sensitivity-based allocation; the
        identification of a monthly feedback loop's measurement problem
        (one ${\sim}95\%$-noise observation per update); and a
        bit-reproducible pipeline --- including a data-determinism defect
        whose fix overturned two previously-drawn conclusions --- with every
        experiment pre-registered, cache-verified, and gated on a
        byte-perfect replication of the prior phase's result.
\end{itemize}

\section{Limitations}
\label{sec:limitations}

Stated without hedging. \textbf{(i)} The decomposition's components are
$0.01$--$0.03$ Sharpe: near the 18-year detection floor, individually
non-significant against honest controls, and plausibly underpowered ---
the point estimates are precise, the inference is thin. \textbf{(ii)} The
orientation effect exists under one discovery method (DYNOTEARS); until
another method reproduces it, it should be read as a property of DYNOTEARS'
contemporaneous structure, not of markets. \textbf{(iii)} Long-only US
large-cap equity only; every variant absorbs the full ${\approx}51\%$ GFC
drawdown; transfer to richer universes is untested. \textbf{(iv)} The
universe approximates the S\&P~100 by market-cap ranking and a fixed
snapshot union, with mild residual survivorship effects. \textbf{(v)} The
asset-block SEM marginalises over the macro drivers, so the structural
shocks' assumed independence is an approximation; and the choice of Q ---
though not the experiment --- was informed by the first phase's results,
which the unified \rsNtrials-trial deflation universe only partly corrects.

\section{Broader impact}
\label{sec:impact}

\paragraph{For investors.} Allocation research is consumed by asset managers
and pension schemes, so its failures are borne --- diffusely, and years
later --- by end savers. The characteristic harm in this literature is not a
malfunctioning system but a convincing backtest of an overfit strategy;
multiplicity-corrected, family-wise-honest evaluation is therefore not
methodological decoration but the investor-protection function of research
method, and this report applies it to its own headline. The same honesty is
a corrective to what might be called \emph{causal-washing}: a ``causal''
label invites over-trust, and the finding here is that a causal method's
value can be modest, window- and method-dependent, and --- at short windows
--- largely attributable to structure that cheap symmetrisation already
captures. The label guarantees nothing; the decomposition, not the
branding, is what carries information.

\paragraph{For the field.} Three of this project's technical choices
generalise into norms. The seed audit has an uncomfortable reading for
published neural-pipeline results: an author free to choose among \rsSeedN\
seeds could have reported $\rsSeedMax$ rather than $\rsSeedMin$ --- an
entire claimed edge manufactured by one undisclosed decision --- so
deterministic constructions, or disclosed seed distributions, should be the
default. The measurement problem generalises to any feedback trading
system: a reward that is ${\sim}95\%$ noise per update cannot teach,
whatever the learning rule sitting on top of it, and checking reward
estimability at the design stage would spare the field many dead ends. And
the determinism episode of Section~\ref{sec:cache} --- two conclusions
overturned by a reproducibility fix --- argues that bit-reproducibility
belongs among publication prerequisites rather than aspirations.

\paragraph{Costs.} The content-keyed cache cut the marginal cost of a
configuration from ${\sim}15$ hours to ${\sim}40$ seconds, which is what
made an \rsNtrials-configuration study --- and hence its multiplicity
correction --- affordable at all. The compute and energy economics of
caching-first experimental design are a transferable, if unglamorous,
contribution to sustainable empirical research: rigour and efficiency were
complements here, not a trade-off.

\section{Future work}
\label{sec:future}

\textbf{(1) A second orientation source.} The decomposition's biggest caveat
is method-dependence: repeating the fixed-graph ablation on graphs from a
different family (e.g.\ the non-linear NTS-NOTEARS, whose prior-transfer was
verified in a reduced-scope probe but whose ${\sim}10\times$ fit cost
deferred a full backtest) would test whether orientation value is a
DYNOTEARS artefact. \textbf{(2) Fix the measurement problem, then revisit
feedback.} The monthly loop failed for a measurement reason, not a learning
one: a higher update frequency would give the learner more reward
observations to average over, and a longer reward window would raise the
signal-to-noise of each one. Only after one of these changes is any feedback
rule worth re-testing. \textbf{(3)}
\textbf{Signed structural allocation.} The long-only ERC on
$\Sigma_{\mathrm{struct}}$ is a projection of exact per-shock parity, which
generally requires shorting; a long--short mandate is the natural home for
structural-shock parity. \textbf{(4) Richer universes}, where cross-asset
causal chains are plausibly longer-lived than among 99 co-moving large-caps
--- multi-asset-class portfolios are where the skeleton/orientation split
could genuinely differ. \textbf{(5) Orientation-aware risk attribution.} The
total-effect matrix $B$ supports ex-post decompositions (which shocks drove
this month's loss?) that the correlation toolkit cannot express ---
worthwhile even where the allocation gain is modest. The reproducibility
infrastructure built here --- the deterministic pipeline, the discovery cache
and the ablation harness --- makes each of these a matter of days, not
months, of compute; which, as this project found twice, is a prerequisite
for trusting the answers.
